\section{The local Hecke operators on the original seed}

Fix the data of B77R2--B77R5. In particular the stage and its valid
column, field embedding, level characters, and the standard self-dual
frames on the terminal row and on each lift row are fixed. The field
extension on the host is unramified at finite places. Write $J$ for
the full level subgroup on whichever row and tuple are under
consideration and $d(y)$ for its decorated lattice at a finite adelic
argument $y$. Each statement below uses the coefficient reading already
chosen on that row.

% BEGIN H78K1
\paragraph{H78K1. Extending the split cosets to the full level.}
Let $G_v$ denote the subgroup supported at a single finite place $v$
lying away from the levels and their conjugates. Its intersection with
$J$ is the integral unitary stabilizer $K_v$. Indeed, in the standard
frame level membership requires that a matrix and its inverse preserve
the integral lattice and that their marked reductions satisfy the flag
conditions. A supported matrix is the identity at all those markings.
Conversely, for any $j\in J$ retain just its $v$-component, putting
identity elsewhere. This is unitary and integral with trivial marked
reductions, hence belongs to $K_v$. The complementary factor belongs
to $J$ and commutes with $G_v$. These assertions persist under the
row's frame isomorphism, since conjugation over the finite adeles
acts component by component.

Let $h\in G_v$. If $K_vhK_v=\coprod_{a\in A}aK_v$ with $A\subset
G_v$, then $JhJ=\coprod_{a\in A}aJ$. In fact write the left factor
of $jhj'$ as $k_v k^v$. Commuting $k^v$ past $h$ puts $jhj'$ in
$k_vhJ$, hence in one $aJ$. Conversely all these cosets belong to
$JhJ$. If $aJ=bJ$ then $a^{-1}b\in J\cap G_v=K_v$, so $a=b$.
Also $\chi$ is one on $K_v$, since each local factor in its marked
determinant character is evaluated on the identity. In B77R5 the
representatives already belong to $G_v$: each lies in the product
$K_vhK_v$, and every factor has support at $v$. The same finite
set thus gives both the full double coset and the prescribed
neighbor bijection for every $y$.
% END H78K1

% BEGIN H78K2
\paragraph{H78K2. The split determinant and the direction of the Hecke move.}
For the source label let $P$ be the prime indexing the neighbor set
and $\bar P$ its conjugate. Choose $\pi\in\mathcal O_E$ with valuation one at
$\bar P$ and zero at $P$. Such an integer is obtained by prescribing a nonzero class in $\bar P/\bar P^2$ and the class of one modulo $P$ by the Chinese remainder theorem. Set $\lambda=\pi/\bar\pi$.
For $1\leq d\leq n$, use the element $\eta$ equal at $v$ to
$\operatorname{diag}(\lambda I_d,I_{n-d})$ and identity off $v$.
Its finite image, in the actual row coordinates, is the operator
representative. The source predicate for that representative calls
for the \emph{inverse} adelic move on the $P$-component.
Extend the finite element by the identity at infinity and denote the
extension, including the frame transport, by $\alpha$. The split
idempotent $e$ selects the $P$-component and its conjugate selects
$\bar P$. Their sum supports $\alpha$ and $\alpha^{-1}$. In the source predicate take the witness $\alpha^{-1}$, whose inverse has finite image equal to the operator representative.
At $P$, $\bar\pi$ has valuation one, $\pi$ is a unit, and
$\alpha^{-1}$ in an integral frame has $d$ diagonal entries
$\bar\pi/\pi$, with the other entries one. Consequently
\[
  e\bar\pi\,\widehat{\Lambda}
    \ \subseteq\ e\alpha^{-1}\widehat{\Lambda}
    \ \subseteq\ e\widehat{\Lambda}.
\]
For the first inclusion the preimage on the first $d$ coordinates
is obtained by multiplication by $\pi$, and on the remaining
coordinates by $\bar\pi$. Both are integral. Integral change of
frame with integral inverse transports both inclusions.

The determinant of $\alpha^{-1}$ on $e$ is
$\bar\pi^{\,d}\pi^{-d}$. The full unit witnessing the unit factor
in the source determinant condition is
$u=e\pi^{-d}+1-e$, with inverse $e\pi^d+1-e$.
Integrality at $P$ follows from $\operatorname{ord}_P(\pi)=0$,
and every other component is one. Determinants commute with
reindexing, scalar extension and conjugation, so this computation
is unchanged on the actual row. This gives the source's split
label, including its determinant witness. H78K1 and B77R4--B77R5
now give the positive-degree split Hecke comparison for both
presentations, with precisely their input prime $P$; the
degree-zero identity and degrees above the rank are B77R4.
% END H78K2

% BEGIN H78K3
\paragraph{H78K3. Why elementary divisors detect the unitary double coset.}
Consider a conjugation-fixed finite prime on the host. The local
coefficient extension is unramified quadratic. Its integer rings
will be written $O/O_0$, with involution $c$ and common uniformizer
$\varpi\in O_0$. A self-dual Hermitian lattice has an integral
orthonormal basis: its reduction is nonsingular Hermitian over
the quadratic residue field and has a vector of nonzero norm.
Lift the vector and rescale its unit norm to one, then split its
orthogonal summand and continue by induction. The norm on units
is onto by the finite-field norm and successive lifting with the
residue trace. Thus one may equally choose the antidiagonal
Hermitian matrix in an integral basis, by applying the same argument
to that matrix. This works in residue characteristic two.

We give the integral factorization needed when two unitary matrices
have the same elementary divisors. First, if $f\in O_0[X]$ is monic
with unit constant term, the class of $X$ in $A_0=O_0[X]/f$ is
a unit and is a norm from $A=A_0\otimes_{O_0}O$. Here are details
allowing repeated factors of the reduction of $f$. Modulo the
Jacobson radical $J$ of $A_0$ the base is a product of finite fields.
On each of them the scalar extension is either quadratic or split.
The norm is surjective in both cases. Lift a residue solution to a
unit of $A$. The trace $O\to O_0$ has an element of trace one,
since the residue trace is onto. A lift of a residue element of trace
one has unit trace, and division by that unit gives trace one. Thus for each $z\in J^m$ one can choose
$b\in J^m A$ with $b+\bar b=z$. Modulo $J^{m+1}$ the norm
of $1+b$ is $1+z$ for $m\ge1$. Successive corrections remove the
norm error. The algebra is complete for this filtration: it is
finite free over the complete DVR and $J^a\subseteq\varpi A_0
\subseteq J$ for some $a$, since the radical modulo $\varpi$
is nilpotent. The limiting solution is represented by a polynomial
$p$ with $p\bar p\equiv X\bmod f$.

Work in an orthonormal basis and let $*$ denote conjugate transpose.
If $g'=hgk$ are unitary $g,g'$ with $h,k\in\mathrm{GL}_n(O)$,
put $B=h^*h$ and $B'=(kk^*)^{-1}$. Then $Bg=gB'$.
Their common characteristic polynomial is monic, has unit constant
term and is fixed by $c$, since $B$ is self-adjoint. Choose the
same polynomial $p$ as above for both matrices. By Cayley--Hamilton
\[
 p(B)^*p(B)=B,\qquad p(B')^*p(B')=B',\qquad p(B)g=g p(B').
\]
The unit determinants on the right make both polynomial matrices
integrally invertible. Therefore
\[
   g'=\bigl(hp(B)^{-1}\bigr)\,g\,\bigl(p(B')k\bigr)
\]
has integral unitary outside factors. Equality of ordinary
elementary divisors gives the initial $h,k$ by Smith normal form
over the DVR after clearing a common power of $\varpi$.
The reverse implication follows from integral invertibility.

If $s_1\ge\cdots\ge s_\ell>0$ and $2\ell\le n$, the length-$n$
list
\[
   (s_1,\ldots,s_\ell,0,\ldots,0,-s_\ell,\ldots,-s_1)
\]
occurs: in the antidiagonal integral frame scale the paired
coordinates by $\varpi^{s_i}$ and $\varpi^{-s_i}$.
This preserves each antidiagonal pairing. Extend by identity at
all other places. The intersection of the resulting conjugate of
$K_v$ with $K_v$ is open and of finite index (the latter group is
compact). Hence the labeled double coset has a finite set of
representatives supported at $v$. This includes the empty list,
where the move may be the identity.
% END H78K3

% BEGIN H78K4
\paragraph{H78K4. Inert neighbors detect the entire label.}
Write $N$ for the starting global lattice and $N'$ for a translated
one. In a pair of integral bases at the inert prime $Q$, an elementary
divisor $a>0$ contributes $O/Q^a$ to $N/(N\cap N')$.
A negative exponent contributes its absolute value to the analogous
quotient of $N'$. Here the local notation $O/Q^a$ also represents
the original global module $\mathcal O_E/Q^a$. Localization gives
the same quotient because it is supported at $Q$.
For a unitary relative matrix the sorted list of divisors is its
own negative reverse: inversion negates and reverses the list,
while conjugate transpose preserves it. Therefore a move with
positive list $s$ gives the two quotient types required of an
inert neighbor. Its marked decorations agree away from $Q$.

Conversely take a decorated unimodular neighbor of type $s$.
The pair of quotient conditions implies that the lattices coincide
at every finite prime other than $Q$. At $Q$ their self-dual
completions admit integral orthonormal frames by H78K3. Mapping
one frame to the other defines a unitary isometry. In the adelic
coordinates of the starting lattice use this isometry there and
identity elsewhere; equality of lattices can be tested after
localization, so the resulting lattice is precisely $N'$. The
off-$Q$ decorations are the original ones by the neighbor condition.
For a finite module $M=\bigoplus_j\mathcal O_E/Q^{a_j}$ one has
\[
 \log_{\#k(Q)}\# M[Q^m]=\sum_j\min(a_j,m).
\]
Successive differences recover the number of exponents at least
$m$, hence the full positive list. Together with the rank and
self-duality this forces exactly the label of H78K3. That claim
now identifies the double coset of the supported isometry. Two
supported elements giving the same neighbor differ on the right by
the integral stabilizer, so the representative is unique. This
proves the bijection simultaneously for all translates of the
starting framed lattice. Integral marking changes on a lift row
preserve the local calculation: in the source notation the increment
is conjugated by $Ak$ for the reading $A^{-1}yAk$.

Use H78K1 at this off-level place to pass to $J$. The level character
on the supported stabilizer is one. For each form $f$ sum
its coefficient values over the neighbor bijection, and apply the
additive coefficient map of B77R2 or B77R3. This proves
\[
  j(T_{Q,s}f)(y)=\sum_{a\in A_{Q,s}} j(f)(ya)
\]
with the ordinary and the lift ranks respectively and with the
original normalization. In the guarded record a label outside the
allowed lengths or orders can be assigned $1$ and the empty set:
the assertions about its coset and operator are required only for
admissible labels. All representatives for an admissible label
are chosen before $f$ and $y$.
% END H78K4

% BEGIN H78K5
\paragraph{H78K5. Right translation and the character in the finite flag action.}
Write $Q_t$ for the parahoric group acting on a tuple's reading,
$\mathrm{red}:Q_t\to G_0$ for its reduction, $\chi_Q$ for its
character, and $m$ for the fixed semilinear coefficient reading.
Use the placewise identification $e_G$ to read the finite group.
The map $\mathrm{red}$ is onto. At each of the marked split pairs
lift the prescribed invertible leading block from the residue field
to an integral invertible matrix (lift the entries, then check its
determinant is a unit). Extend by the identity on the other diagonal
blocks. At the conjugate prime use the inverse adjoint in the
self-dual frame. Take identity at the disjoint auxiliary markings
and elsewhere. This is an adelic parahoric element with the
specified tuple of reductions. The finite lift of the element works
on the given tuple on the ordinary row and on every lift row.

For $\kappa\in Q_t$, put $g=\mathrm{red}(\kappa)$. In the finite
decorated model let $z$ be the full residue change of frame and
$b(g)$ the prescribed leading-block action. Right translation
changes the decoration by the inverse residue-frame action, since
the decoration is read using the inverse coordinate matrix.
The correction in the stabilizer is $h=z^{-1}b(g)$.
At a type place $z$ is block triangular and its first block is the
chosen reduction, so the corresponding block of $h$ is the identity.
The later diagonal blocks of $h$ are the inverses of those of $z$.
At an auxiliary tuple place $b(g)=1$ and $z$ preserves its two-block
flag. These are precisely the type-level conditions on $h$.
On each block the determinant of the inverse is the inverse of the
determinant. Substitute these determinants in the character of the
inverse type, then apply the given character embedding. On a lift
row the factors at the type places are $\chi_{i+1}^{-1}$ on the
prescribed secondary block determinant, and the tuple places
contribute the prescribed $\xi_t$ on their complementary block
determinants. On the ordinary row use all its successive diagonal
blocks in the same determinant product. Multiplicativity and
disjointness of the markings give
$\sigma(\chi_{\rm type}(h))=\chi_Q(\kappa)$ in the inverse-type
convention. This calculation uses the character identification
of the given column, including $e_G$ on the leading blocks.

Apply equivariance of the decorated form at $h$ and the semilinearity
of $m$. With the source's action $g\cdot f$ the result is
\[
       j(g\cdot f)(y)=\chi_Q(\kappa)^{-1}j(f)(y\kappa).
\]
Indeed $hb(g)^{-1}=z^{-1}$ carries the original marked point to
the translated one; moving $h$ through the coefficient value gives
the indicated inverse character. Each side uses the same $m$ for
every twist of that row. This proves the finite action and its
reduction clause on the presentations of B77R2--B77R3.
% END H78K5

% BEGIN H78K6
\paragraph{H78K6. The Hecke and reading data for the preliminary seed.}
Take the standard seed frames of B77R1--B77R3, their coefficient
models and pairings, the original level and placewise finite group
identifications, the specified archimedean family and terminal
discrete-series model. Then the seed bridge $\mathrm{BrgL}$ and both
standard readings are obtained together. In particular the full-level
split cosets, inert operators and finite-group action of B77R4 are
now supplied on its own ordinary and lift spaces.

Use the terminal pairing of B77R1 and the ordinary and lift
presentations of B77R2--B77R3. On each of these use H78K1--H78K2
for the split test, H78K3--H78K4 for its inert test and H78K5 for
its finite action. Degree zero and degrees beyond the rank are as in
B77R4. For the remaining place maps take the injections of split and
inert prime labels into the primes of the host. The excluded set on
a tuple is the union of the type markings and the twist markings;
the bad set in the real field is the set of contractions of the type
markings. At a split prime the two exclusion tests therefore agree
in both directions. At a conjugation-fixed prime the conjugate
ideal is itself. Conversely the host's unramified CM condition
supplies an element with nontrivial residue conjugation there,
giving the required inert label for every such unexcluded prime.
This verifies the two place comparisons and their surjectivity
statements on their specified domains.

At an allowed split prime the tuple character is unramified: any
ramification would put the prime or its conjugate among the twist
markings, contrary to the exclusion test. Multiply the embedded
base value by the unramified tuple value to obtain the required
$\mu_0\xi_t$ value. The eta \emph{character} value at the base tuple
is the given character pin. The residue roots at the type markings
have the two specified prime orders by the norm condition; those at
the twist markings come from $\mathrm{ValsIn}$ (F77W1 gives its
simultaneous realization). This supplies the remaining arithmetic
fields of the bridge. Finally the stage ranks identify the two
ranks with $r+mt$ and $r+t$, respectively. Using these equalities
in the packet leaves the evaluation maps exactly the ones already
chosen. This supplies the reading input in V77Q2; the weak theta
kernel construction and its analytic inputs, and the production
of such a seed throughout the family, are separate requirements.
% END H78K6
